\section*{अध्याय \ref{ch:g9:climate-past-and-present} --- पृथ्वी की जलवायु, अतीत और वर्तमान}

\begin{solution}{exo:g9:climate-past-and-present:1}
(क) लगभग \qty{-19}{\degreeCelsius}। (ख) लगभग \qty{15}{\degreeCelsius}।
(ग) $15 - (-19) = 34$: लगभग \qty{34}{\degreeCelsius}। (घ) इनमें से दो:
कार्बन डाइऑक्साइड, जलवाष्प, मीथेन।
\end{solution}

\begin{solution}{exo:g9:climate-past-and-present:2}
(क) \omterm{def:g1:watching-the-weather:weather}{मौसम}। (ख) \omterm{def:g5:weather-and-climate:climate}{जलवायु}। (ग) \omterm{def:g5:weather-and-climate:climate}{जलवायु} (और \omterm{def:g9:climate-past-and-present:climate-change}{जलवायु परिवर्तन}: दशकों में होने वाला बदलाव)।
(घ) \omterm{def:g1:watching-the-weather:weather}{मौसम}।
\end{solution}

\begin{solution}{exo:g9:climate-past-and-present:3}
(क) \qty{19}{cm}, ग्यारह वर्ष। (ख) $19 \div 11 \approx \qty{1.7}{cm}$।
(ग) हर वर्ष का \omterm{def:g4:clouds-rain-snow:precipitation}{हिम} पिछले वर्षों के \omterm{def:g4:clouds-rain-snow:precipitation}{हिम} के ऊपर गिरता है, इसलिए परत जितनी गहरी, वह
उतनी पहले गिरी।
\end{solution}

\begin{solution}{exo:g9:climate-past-and-present:4}
(क) $22.4 - 21.0 = +\qty{1.4}{\degreeCelsius}$। (ख) $20.3 - 21.0 =
\qty{-0.7}{\degreeCelsius}$। (ग) \qty{0}{\degreeCelsius}। (घ) वह महीना
अपने 1991--2020 के औसत से ठंडा था।
\end{solution}

\begin{solution}{exo:g9:climate-past-and-present:5}
(क) 1880 के दशक से 1930 के दशक तक, 1950 का दशक और 1960 का दशक। (ख) लगभग
\qty{0.25}{\degreeCelsius} और \qty{0.80}{\degreeCelsius}। (ग) $0.80 -
0.25 = \qty{0.55}{\degreeCelsius}$। (घ) नहीं: 1880 के दशक से 1910 के दशक तक के दशक थोड़े
ठंडे होते गए, और 1950 का दशक 1940 के दशक से थोड़ा ठंडा था; पर 1970 के दशक से हर दशक
अपने से पहले वाले दशक से गरम रहा
है।
\end{solution}

\begin{solution}{exo:g9:climate-past-and-present:6}
(क) $(324.62 - 315.98) \div 10 \approx \qty{0.86}{ppm}$ प्रति वर्ष। (ख)
$(427.35 - 401.01) \div 10 \approx \qty{2.63}{ppm}$ प्रति वर्ष। (ग) $2.63
\div 0.86 \approx 3.1$: लगभग तीन गुना तेज़। (घ) $(427.35 - 315.98)
\div 315.98 \times 100 \approx 35~\%$।
\end{solution}

\begin{solution}{exo:g9:climate-past-and-present:7}
(क) $8 \times 10^5$ वर्ष। (ख) लगभग \num{120000} से \num{130000} वर्ष
पहले: वहाँ तापमान का वक्र 0 से ऊपर उठता है। (ग) लगभग \qty{275}{ppm}।
(घ) हाँ: दोनों वक्रों के शिखर और गर्त एक ही समय पर आते हैं।
\end{solution}

\begin{solution}{exo:g9:climate-past-and-present:8}
(क) $3.2 \times 10 = \qty{32}{mm} = \qty{3.2}{cm}$। (ख) $2100 - 2025 =
75$ वर्ष; $75 \times 3.2 = \qty{240}{mm}$, लगभग \qty{24}{cm} (और
ज़्यादा, अगर बढ़त तेज़ हो जाए)। (ग) \omterm{def:g2:water-on-earth:ocean}{समुद्र} का पानी गरम होने पर फैलता है; ज़मीन पर की
बर्फ़ पिघलती है और उसका पानी \omterm{def:g2:water-on-earth:ocean}{समुद्र} में बह जाता है। (घ) नहीं: \omterm{def:g8:ice-on-earth:sea-ice}{समुद्री बर्फ़} पहले से ही
\omterm{def:g2:water-on-earth:ocean}{समुद्र} पर तैरती है, इसलिए उसका पिघलना \omterm{def:g2:water-on-earth:ocean}{समुद्र} को ऊपर नहीं उठाता।
\end{solution}

\begin{solution}{exo:g9:climate-past-and-present:9}
(क) 90। (ख) 66। (ग) 99। (घ) \emph{संभावित} का अर्थ है 100 में कम से कम 66
संभावनाएँ, निश्चितता नहीं: अख़बार कथन को वैज्ञानिकों के कहे से ज़्यादा पक्का बना
देता है।
\end{solution}

\begin{solution}{exo:g9:climate-past-and-present:10}
एक जगह का एक ठंडा जाड़ा \emph{\omterm{def:g1:watching-the-weather:weather}{मौसम}} है: वह बताता है कि एक जगह पर कुछ महीनों में क्या
हुआ। \omterm{def:g9:climate-past-and-present:global-warming}{वैश्विक तापन} \emph{\omterm{def:g5:weather-and-climate:climate}{जलवायु}} का बदलाव है, दशकों में मापा गया, और \emph{वैश्विक}
है: वह पूरी पृथ्वी का औसत है। यहाँ एक ठंडा जाड़ा तब भी पड़ सकता है जब पूरी पृथ्वी पहले से
गरम हो, और दशकों के औसत बढ़ते ही
जाते
हैं।
\end{solution}

\begin{solution}{exo:g9:climate-past-and-present:11}
(क) $280 - 190 = \qty{90}{ppm}$, लगभग $2 \times 10^4$ वर्षों में। (ख) $427
- 280 = \qty{147}{ppm}$, $2025 - 1750 = 275$ वर्षों में। (ग) दूसरी
बढ़त बड़ी है (147 बनाम 90~ppm)। दरें: $90 \div (2 \times 10^4) =
4.5 \times 10^{-3}$~ppm प्रति वर्ष, बनाम $147 \div 275 \approx 5.3 \times
10^{-1}$~ppm प्रति वर्ष: लगभग 120 गुना तेज़। (घ) यह कि 2019 में वह
कम से कम बीस लाख वर्षों में किसी भी समय से ज़्यादा थी (उच्च विश्वास)।
\end{solution}

\begin{solution}{exo:g9:climate-past-and-present:12}
(क) वे लगभग बराबर हैं: 1.07 बनाम \qty{1.06}{\degreeCelsius}। सबसे अच्छा
अनुमान कहता है कि मानव प्रभाव लगभग पूरे देखे गए तापन को समझाता है। (ख) नहीं:
प्राकृतिक कारणों ने मिलकर तापमान को अधिक से अधिक \qty{0.1}{\degreeCelsius} बदला,
देखे गए तापन से दस गुना कम। (ग) 100 में कम से कम 66 संभावनाएँ कि मानव-जनित
तापन 0.8 और \qty{1.3}{\degreeCelsius} के बीच है।
\end{solution}

\begin{solution}{pb:g9:climate-past-and-present:1}
\textbf{1.} क्योंकि वही संदर्भ अवधि है: विसंगतियाँ उसके अपने औसत से मापी जाती
हैं।

\textbf{2.} ठंडा, \qty{0.22}{\degreeCelsius} से।

\textbf{3.} गरम, \qty{0.98}{\degreeCelsius} से।

\textbf{4.} $0.98 - (-0.22) = \qty{1.20}{\degreeCelsius}$।

\textbf{5.} $427 - 316 = \qty{111}{ppm}$।

\textbf{6.} $111 \div 316 \times 100 \approx 35~\%$।

\textbf{7.} $427 - 300 = \qty{127}{ppm}$।

\textbf{8.} कार्बन डाइऑक्साइड गरम सतह से निकलने वाले अवरक्त प्रकाश का एक हिस्सा सोखती
है और उसका एक भाग वापस नीचे भेजती है। हवा में उसके ज़्यादा होने पर ज़्यादा ऊर्जा सतह पर
लौटती है, जो गरम होती है।

\textbf{9.} $2025 - 1993 = 32$ वर्ष; $32 \times 3.2 \approx
\qty{102}{mm}$, लगभग \qty{10}{cm}।

\textbf{10.} गरम होता समुद्री पानी फैलता है; ज़मीन पर की बर्फ़ (\omterm{def:g8:ice-on-earth:glacier}{हिमनद}, \omterm{def:g8:ice-on-earth:ice-sheet}{हिमचादरें})
पिघलकर \omterm{def:g2:water-on-earth:ocean}{समुद्र} में जाती है।

\textbf{11.} हाँ, उतने ही आकार का: लगभग \qty{1.2}{\degreeCelsius} बनाम
\qty{1.09}{\degreeCelsius}। वे थोड़े अलग हैं क्योंकि अवधियाँ एक जैसी नहीं हैं
(1850--1900 की जगह 1880--1899; 2011--2020 की जगह 2015--2024: ज़्यादा हाल के वर्ष
ज़्यादा गरम थे), और क्योंकि IPCC कई अभिलेखों को मिलाता है, केवल NASA के अभिलेख को
नहीं।

\textbf{12.} ``उन्नीसवीं सदी के अंत से पृथ्वी की सतह लगभग
\textbf{\qty{1.2}{\degreeCelsius}} गरम हुई है।''
\end{solution}
